% ch4.tex
\chapter{আসিরিয়ান জ্যোতিষশাস্ত্র}

\epigraph{``\textit{বল মহাবিশ্বের মহাকাশ ফাড়ি’\\
চন্দ্র সূর্য গ্রহ তারা ছাড়ি’\\
ভূলোক দ্যুলোক গোলক ভেদিয়া\\
খোদার আসন ‘আরশ’ ছেদিয়া,\\
উঠিয়াছি চির-বিস্ময় আমি বিশ্ববিধাতৃর}''}{\textit{নজরুল}}


খ্রিস্টপূর্ব প্রথম সহস্রাব্দের শুরুতে উত্তরে, টাইগ্রিসের উজানে, একটি নতুন সেমিটিক শক্তির উত্থান ঘটে। এর রাজধানীর নাম ছিল আশুর, যা প্রধান দেবতারও নাম। ব্যাবিলনসহ আশপাশের রাষ্ট্রগুলোর সঙ্গে অবিরাম সংঘাতের মধ্য দিয়ে আশুর বিস্তার লাভ করে এবং মেসোপটেমীয় সমভূমির শাসকশক্তিতে পরিণত হয়। খ্রিস্টপূর্ব ৮০০ সালের দিকে আসিরিয়া নিকট-এশিয়ার সবচেয়ে শক্তিশালী রাষ্ট্র হয়ে ওঠে। এর রাজারা ছিলেন তিগলাথ-পিলেসার (খ্রিস্টপূর্ব ৭৪৫–৭২৭), শালমানেসার (৭২৬–৭২২), সারগন (৭২২–৭০৫), সেনাখেরিব (৭০৫–৬৮২), এসারহাদ্দোন (৬৮২–৬৬৮) এবং আশুরবানিপাল (খ্রিস্টপূর্ব ৬৬৮–৬২৬)। তাঁদের আমলে একের পর এক বড় এবং প্রায়ই নিষ্ঠুর যুদ্ধে সিরিয়া, প্যালেস্টাইন, ফিনিশিয়া, এমনকি কখনো কখনো মিশরও জয় করা হয়। সীমানা এশিয়া মাইনর, আর্মেনিয়া ও মিডিয়ার দিকে বিস্তৃত হয়। মার্কিন অ্যাসিরিওলজিস্ট ওলমস্টেডের মতে, আসিরীয় বিজয়ের প্রধান বস্তুগত কারণ ছিল অস্ত্রে লোহার ব্যবহার।

নতুন রাজধানী নিনেভে ছিল একটি বিশাল সামরিক সাম্রাজ্যের রাজনৈতিক কেন্দ্র, আর সে হিসেবে একে অপূর্ব সব ভবন দিয়ে সাজানো হয়েছিল। কিন্তু ব্যাবিলন ছিল একটি বড় ও সমৃদ্ধ বাণিজ্যকেন্দ্র, যার ধনী নাগরিকেরা অনেকাংশে নিজেরাই নিজেদের শাসন করত। প্রাচীন সংস্কৃতির পূজনীয় পীঠস্থান হিসেবে ব্যাবিলন তার মর্যাদা ধরে রেখেছিল। আসিরীয় রাজারা এর গুরুত্ব স্বীকার করতেন। তাঁরা নিজেরা ব্যাবিলনে গিয়ে "মারদুকের হাত ধরতেন", অর্থাৎ আনুষ্ঠানিকভাবে শাসনভার গ্রহণ করতেন, অথবা কোনো আত্মীয়কে অধীনস্থ রাজা হিসেবে মনোনীত করতেন।

রুক্ষ ও যুদ্ধপ্রিয় বিজেতারা প্রায়ই যা করে, আসিরীয়রাও তা-ই করল: তারা বিজিতদের সংস্কৃতি গ্রহণ করল এবং তার রূপ ও ঐতিহ্য বহন করে চলল। আসিরীয় দেবমণ্ডলী ব্যাবিলনেরটির মতোই ছিল, পার্থক্য কেবল এই যে এখন আশুর হলেন দেবতাদের নতুন দেবতা। একই পুরোহিতসমাজ নতুন শাসকদের সেবায় পুরোনো পবিত্র রীতিতে চিন্তা করত, কাজ করত ও লিখত, এবং আগের মতোই একই সামাজিক দায়িত্ব পালন করত। ভাস্কর্যকলার নতুন করে বিকাশ ঘটল। এর এক কারণ, কারিগরেরা ছিল নতুন প্রাণশক্তিসম্পন্ন একটি নতুন জাতি, আর নতুন শাসকেরা বিজয় ও লুণ্ঠনের মাধ্যমে ধনী হয়েছিলেন। আরেক কারণ, নিম্নভূমির ইটের বদলে আশপাশের পাহাড়ের পাথর নির্মাণসামগ্রী হিসেবে ব্যবহৃত হতো। সুন্দর উদ্‌গত খোদাইয়ের জন্য এটি ছিল চমৎকার উপকরণ, আর সেই খোদাইগুলো আজ পশ্চিমের জাদুঘরগুলোর শোভা বাড়াচ্ছে। এই বিকাশ চূড়ায় পৌঁছায় যখন সম্পদ ও সংস্কৃতি বাড়ার সঙ্গে সঙ্গে প্রথম দিকের কঠোর যোদ্ধাদের জায়গায় এলেন এমন রাজপুত্রেরা, যাঁরা শিল্প ও বিজ্ঞান ভালোবাসতেন এবং তার পৃষ্ঠপোষকতা করতেন। তখনই আশুরবানিপাল তাঁর প্রাসাদে একটি গ্রন্থাগার স্থাপন করেন। তিনি ব্যাবিলোনিয়ার সব প্রাচীন স্থান ও মন্দির থেকে পুরোনো লেখাগুলো সংগ্রহ করে অনুলিপি করার আদেশ দেন। হাজার হাজার মাটির ফলক সুশৃঙ্খল সারিতে রাখা হয়েছিল, পাশে তাদের শিরোনাম লেখা থাকত। ক্যাটালগ, অভিধান ও ভাষ্য দিয়ে এগুলোর পরিপূরক ও ব্যাখ্যা দেওয়া হতো। নতুন নতুন মহাফেজখানার দলিল, রাজার কাছে পাঠানো প্রতিবেদন এবং কর্মকর্তাদের সঙ্গে তাঁর চিঠিপত্র দিয়ে সংগ্রহটি অবিরাম বাড়ানো হতো। কেবল এই স্থান থেকেই খুঁড়ে পাওয়া ১৩,০০০-এর বেশি খণ্ড ব্রিটিশ মিউজিয়ামে সংরক্ষিত আছে, আর অন্যান্য স্থান থেকে আরও বহু হাজার পাওয়া গেছে। এগুলো থেকে সেই সমাজের রীতিনীতি ও ধ্যানধারণা, ব্যবসায়িক জীবন ও সংস্কৃতি, ধর্ম ও অর্থনীতি, এবং তার জ্যোতির্বিজ্ঞানেরও একটি ভালো চিত্র পাওয়া যায়।

অসংখ্য লেখায় আগের যুগের পঞ্জিকাই পাওয়া যায়। চান্দ্র মাস শুরু হতো বাঁকা চাঁদের প্রথম আবির্ভাবের সন্ধ্যা থেকে, আর ১২টি মাসের সঙ্গে মাঝে মাঝে একটি ত্রয়োদশ মাস যোগ করে বছর পূর্ণ করা হতো। স্পষ্ট উল্লেখ না থাকলেও এটি প্রায় নিশ্চিত যে নক্ষত্রসংক্রান্ত ঘটনা দিয়ে, সাধারণত নক্ষত্রের প্রভাতী উদয় দিয়ে, অধিমাস যোগ নিয়ন্ত্রণ করা হতো। আগে উল্লিখিত ৩ × ১২টি মাস-নক্ষত্রের তালিকার একটি অনুলিপি আশুরবানিপালের গ্রন্থাগারে আছে, যা এর একটি ইঙ্গিত হতে পারে। জর্জ স্মিথের প্রকাশিত একটি লেখায় আরেকটি পদ্ধতির ইঙ্গিত পাওয়া যায়: "নিসান্নুর প্রথম দিনে চাঁদ ও মুলমুল নক্ষত্র [কৃত্তিকা] একসঙ্গে থাকলে বছরটি সাধারণ; নিসান্নুর তৃতীয় দিনে চাঁদ ও মুলমুল নক্ষত্র একসঙ্গে থাকলে বছরটি পূর্ণ।" বাক্যের শেষ অংশের অর্থ হলো, কৃত্তিকা সূর্যাস্তের অনেক পরেও দৃশ্যমান থাকে। বসন্তের এত আগে এমন হলে একটি ত্রয়োদশ মাস যোগ করা প্রয়োজন।

নক্ষত্র ও তারামণ্ডলগুলোকে তিনটি ভাগে ভাগ করা হয়েছিল: উত্তরের, মধ্যের ও দক্ষিণের তারামণ্ডল। এগুলো যথাক্রমে এনলিল, আনু ও এয়ার এলাকা নামে পরিচিত ছিল।

তবে নক্ষত্র পর্যবেক্ষণের প্রধান উদ্দেশ্য আর পঞ্জিকা ছিল না। আসিরীয় যুগে মানুষের চিন্তা ও চর্চাকে সবচেয়ে জোরালোভাবে নির্ধারণ করত জ্যোতিষশাস্ত্র, অর্থাৎ এই ধারণা যে নক্ষত্রের গতিপথ পৃথিবীর ঘটনাবলির জন্য তাৎপর্যপূর্ণ। এখন আকাশের ঘটনাবলি অনেক গভীর আগ্রহ নিয়ে অধ্যয়ন করা হতো মানুষের ভাগ্য সম্পর্কে, বিশেষ করে রাজা ও সাম্রাজ্যের ভাগ্য সম্পর্কে, শুভাশুভ লক্ষণ খুঁজতে।

শুভাশুভ লক্ষণে এই বিশ্বাস আদিম মানুষের মধ্যেও ছিল। এটি ছিল তার এই বিশ্বাসের স্বাভাবিক পরিণতি যে সে অদৃশ্য আত্মিক শক্তি দিয়ে পরিবেষ্টিত, আর সেই শক্তিগুলো তার জীবন ও কাজকে প্রভাবিত করে। তাদের অনুগ্রহ ও সাহায্য লাভ করা, তাদের শত্রুতা প্রশমিত বা প্রতিহত করা এবং তাদের অভিপ্রায় জানা ছিল তার কাছে জীবন-মরণের প্রশ্ন। ভূত তাড়ানো, নৈবেদ্য, মন্ত্র, তাবিজ এবং কাজের সঙ্গে জড়িত জাদু তার দৈনন্দিন জীবন জুড়ে থাকত। এই আত্মাদের বেশিরভাগের বাস ছিল আকাশে। তাই প্রাচীন কালেই মেসোপটেমীয় পুরোহিতদের মনে নক্ষত্র ও মানুষের ভাগ্যের মধ্যে ঘনিষ্ঠ সম্পর্কের ধারণা জন্মেছিল। স্বর্গ খুব দূরে ছিল না; পৃথিবীর শাসক ও তাদের প্রজারা ওপরের দেবতাদের কাছাকাছিই ছিল। লাগাশের রাজা গুদেয়া (আনুমানিক খ্রিস্টপূর্ব ২৫০০) একটি পাথরের সিলিন্ডারে খোদাই করা লেখায় একটি মন্দির নির্মাণের বর্ণনা দিয়েছেন। তিনি বলেছেন, কীভাবে স্বপ্নে এয়ার কন্যা দেবী নিসাবা তাঁর সামনে আবির্ভূত হয়েছিলেন: "তাঁর হাতে ছিল উজ্জ্বল লেখনী; তিনি শুভ আকাশ-লক্ষণযুক্ত একটি ফলক বহন করছিলেন এবং চিন্তামগ্ন ছিলেন।" আরও পরে: "তিনি মন্দির নির্মাণের জন্য শুভ নক্ষত্রের কথা ঘোষণা করলেন।" প্রথম ব্যাবিলনীয় সাম্রাজ্যের নিন-দার-আন্না লেখায় গ্রহসংক্রান্ত প্রতিটি ঘটনার পর একটি করে লক্ষণ পাওয়া যায়। এতে কোনো সন্দেহ নেই যে এই ভবিষ্যদ্বাণীগুলো জোগানোর জন্যই জ্যোতির্বৈজ্ঞানিক তথ্য সংগ্রহ করা হতো।

ব্যাবিলন ও অন্যান্য নগরের মন্দিরের পুরোহিতদের মধ্যে এমন বিশ্বাস ছিল একটি জীবন্ত শক্তি। ব্যাবিলনীয় সংস্কৃতির প্রভাবাধীন সব অঞ্চলে, আসিরিয়াসহ, এটি ছড়িয়ে পড়ে। বিজয়ের মাধ্যমে সাম্রাজ্য বিস্তারে উদ্‌গ্রীব শক্তিশালী রাজাদের দরবারে ভবিষ্যৎ আগে থেকে জানার প্রয়োজন ছিল বিশেষভাবে বেশি, আর সেখানে জ্যোতিষশাস্ত্র উর্বর ক্ষেত্র খুঁজে পেল। রাজদরবারের জ্যোতিষীদের প্রতিটি বড় উদ্যোগের জন্য লক্ষণ খুঁজে বের করতে হতো। সব গুরুত্বপূর্ণ মন্দির থেকে রাজা নিয়মিত প্রতিবেদন পেতেন যে আকাশে কী ঘটছে এবং তার ব্যাখ্যা কী। এই প্রতিবেদনগুলো সেই মহাফেজখানায় সংরক্ষিত ছিল যাকে আমরা আশুরবানিপালের গ্রন্থাগার বলি, আর ব্যাখ্যার সুবিধার জন্য সব পুরোনো তথ্যের অনুলিপি সেখানে জড়ো করা হয়েছিল।

অবশ্য সব লক্ষণ নক্ষত্র থেকে আসত না। লক্ষণ সর্বত্রই পাওয়া যেত: বলির পশুর যকৃতে, যা পরিবর্তনশীল বলে লক্ষণের প্রচুর উৎস ছিল; পাখির ওড়ায়; গর্ভপাত, ভূমিকম্প, মেঘ, রংধনু ও জ্যোতির্বলয়ে। এমন সব নির্দেশিকা ছিল, যাতে প্রতিটি ঘটনা লিপিবদ্ধ থাকত, আর সেই সঙ্গে থাকত কিংবদন্তির রাজা আগাদের সারগনের কোন কাজ বা অভিজ্ঞতার সঙ্গে এর কোনোভাবে সম্পর্ক আছে।

স্বাভাবিকভাবেই এগুলোর মধ্যে জ্যোতির্বৈজ্ঞানিক ঘটনা বড় জায়গা জুড়ে ছিল। যেসব জ্যোতিষ্কের ঘটনায় বৈচিত্র্য ও অনিয়ম বেশি, সেগুলোই পূর্বাভাসের জন্য সবচেয়ে উপযোগী ছিল। তাই সবচেয়ে গুরুত্বপূর্ণ ছিল সূর্য ও স্থির নক্ষত্র নয়, বরং চাঁদ ও গ্রহগুলো। স্থির নক্ষত্রগুলো ঘটনাবলির জন্য অপরিবর্তনীয় পটভূমি হিসেবে কাজ করত। ন্যায়বিচারের সর্বদ্রষ্টা রক্ষক সূর্যদেব শামাশ দিনের পর দিন আকাশজুড়ে তাঁর অপরিবর্তনীয় পথে চলেন। বড়জোর তিনি স্বাভাবিকের চেয়ে নিষ্প্রভ বা বেশি লাল হতে পারতেন, কিংবা মেঘ বা গ্রহণের কারণে অদৃশ্য হতে পারতেন। অন্যদিকে চন্দ্রদেব সিন বারবার ফিরে আসা নানা রকম ঘটনার প্রাচুর্য দেখান। গ্রহগুলোর ক্ষেত্রেও একই কথা খাটে। তারা নক্ষত্রদের মাঝে সবচেয়ে অপ্রত্যাশিত উপায়ে ঘুরে বেড়ায়। কখনো স্থির দাঁড়িয়ে থাকে বা উল্টো পথে চলে, কখনো নিজেদের মধ্যে বা উজ্জ্বল নক্ষত্রগুলোর সঙ্গে মিলে সদা পরিবর্তনশীল বিন্যাস তৈরি করে। এদের মনে হতো তারাভরা আকাশের প্রান্তরে স্বেচ্ছায় ঘুরে বেড়ানো জীবন্ত সত্তা, আর ক্রমশ এরাই ব্যাবিলনীয় পুরোহিতদের মনোযোগের প্রধান বিষয় হয়ে উঠল। এরা ছিল সেই মহান দেবতাদের নক্ষত্র, যাঁরা জগৎ শাসন করতেন এবং এই উজ্জ্বল জ্যোতিষ্কগুলোর মধ্য দিয়ে নিজেদের প্রকাশ করতেন।

গ্রহগুলোর নামের আগে বসানো নির্ধারক চিহ্নটি সাধারণত ছিল "নক্ষত্র"-এর চিহ্ন, এখন প্রায়ই তা হয়ে গেল "দেবতা"-র চিহ্ন। শুক্র (যার নাম দিল-বাত) ছিল ইশতারের নক্ষত্র। বৃহস্পতি (উমানপা-উদ্দা, পরে বেশিরভাগ সময় সাগমেগার) ছিল মারদুকের নক্ষত্র। শনি ছিল নিনিবের নক্ষত্র, যিনিও একজন সৌর দেবতা। আর মঙ্গল ছিল মহামারির দেবতা নের্গালের নক্ষত্র। লাল মঙ্গলকে অশুভ নক্ষত্র মনে করা হতো, বৃহস্পতিকে শুভ নক্ষত্র। কিন্তু কার জন্য শুভ বা অশুভ? তা নির্ভর করত সময় ও স্থানের ওপর।

চারটি দেশের প্রতিটির জন্য আলাদা আলাদা মাস, আলাদা দিক এবং আলাদা তারামণ্ডল নির্দিষ্ট ছিল: আক্কাদ (ব্যাবিলোনিয়া), এলাম (পূর্বের পাহাড়ি অঞ্চল), আমুররু (পশ্চিমের মরুভূমি, পরে সিরিয়া) এবং সুবারতু (উত্তর)। একজন আসিরীয় জ্যোতিষী বলেছিলেন, "আমরাই সুবারতু।" বিভিন্ন তারামণ্ডলে গ্রহগুলোর আবির্ভাব, তাদের অবস্থানের স্থায়িত্ব বা দ্রুত গতি, একে অপরের সঙ্গে বা চাঁদের সঙ্গে তাদের মিলন ঘটনার একেবারে অন্তহীন বৈচিত্র্য তৈরি করত। এতে জ্যোতিষীরা ঐতিহ্য ও উদ্ভাবনী বুদ্ধি খাটানোর বিস্তৃত সুযোগ পেতেন। ব্যক্তিগত সিদ্ধান্ত টানার সময় মেলানোর ক্ষমতা ও কল্পনা কাজে লাগানোর, আবার বিচক্ষণভাবে চাটুকারিতা করারও সুযোগ থাকত।

দৃষ্টান্ত হিসেবে এখানে কয়েকটি লেখা দেওয়া হলো। এগুলো অ্যাসিরিওলজির শিক্ষার্থীদের জন্য আর. সি. টমসনের প্রকাশিত সংকলন থেকে নেওয়া।

মারদুকের নক্ষত্র বছরের শুরুতে দেখা দিলে সেই বছর শস্য সমৃদ্ধ হবে। বুধ নিসান্নুতে দেখা দিয়েছে। কোনো গ্রহ [বুধ] লি নক্ষত্রের [অ্যালডেবারান] কাছে এলে এলামের রাজা মারা যাবেন।… বুধ বৃষ রাশিতে দেখা দিয়েছে; এটি শুগি [কৃত্তিকা, পারসিউস] পর্যন্ত নেমে এসেছিল [?]।

শুক্র পশ্চিমে অদৃশ্য হয়। আবু মাসে শুক্র নিষ্প্রভ হয়ে অদৃশ্য হলে এলামে হত্যাযজ্ঞ হবে। আবু মাসে প্রথম থেকে ত্রিশতম দিন পর্যন্ত শুক্র দেখা দিলে বৃষ্টি হবে, আর দেশের ফসল সমৃদ্ধ হবে। মাসের মাঝামাঝি শুক্র পূর্বে সিংহ রাশিতে দেখা দিয়েছে। নিরগাল-ইতির থেকে।

শুক্র যখন তার অবস্থান স্থির করে, তখন রাজার আয়ু দীর্ঘ হবে, দেশে ন্যায়বিচার থাকবে। শুক্র যখন এয়ার পথে থাকে…

দুজু মাসে মঙ্গল দৃশ্যমান; এটি নিষ্প্রভ।… মঙ্গল অস্পষ্টভাবে [?] মধ্যগগনে পৌঁছে উজ্জ্বল হয়ে উঠলে এলামের রাজা মারা যাবেন। দেবতা নের্গাল অদৃশ্য হওয়ার সময় যখন ছোট হয়ে আসেন এবং আকাশের নক্ষত্রদের মতো খুব অস্পষ্ট হন, তখন তিনি আক্কাদের প্রতি দয়া দেখাবেন।… মঙ্গল নিষ্প্রভ হলে শুভ; উজ্জ্বল হলে অশুভ। মঙ্গল বৃহস্পতিকে অনুসরণ করলে সেই বছর শুভ হবে। বুল্লুতু থেকে। 
মঙ্গল আল্লুলের [কর্কট] এলাকায় প্রবেশ করেছে। এটি লক্ষণ হিসেবে গণ্য হয় না। এটি থামেনি, অপেক্ষা করেনি, বিশ্রাম নেয়নি; দ্রুত এগিয়ে গেছে। বিল-না-সির থেকে।

বৃহস্পতি পশ্চিমে চলে গেলে নিরাপদ বসবাস হবে, দেশে কোমল শান্তি নেমে আসবে। এটি আল্লুলের সামনে দেখা দিয়েছে। বৃহস্পতি বেলের পথে উজ্জ্বলতা ধারণ করে নিবিরু (হয়ে উঠলে?), আক্কাদ প্রাচুর্যে উপচে পড়বে, আক্কাদের রাজা শক্তিশালী হবেন।… আগুনের মতো একটি বড় নক্ষত্র পূর্ব থেকে উঠে পশ্চিমে অদৃশ্য হলে যুদ্ধে শত্রুসৈন্য নিহত হবে। রাজত্বের শুরুতে বৃহস্পতিকে সঠিক অবস্থানে দেখা গেছে; দেবতাদের প্রভু তোমাকে সুখী করুন এবং তোমার আয়ু দীর্ঘ করুন। দামকার পুত্র আশারিদু থেকে।

এখানে উল্লিখিত অগ্নিময় নক্ষত্রটি সম্ভবত একটি উল্কা ছিল। পরের লেখাটিতেও একই ব্যাপার:

রাতের এক কাস-বু [দুই ঘণ্টা] পার হওয়ার পর একটি বড় নক্ষত্র উত্তর থেকে দক্ষিণে জ্বলে উঠল। এর লক্ষণ রাজার আকাঙ্ক্ষার অনুকূল। আক্কাদের রাজা তাঁর উদ্দেশ্য সফল করবেন। রাজার সেবক জ্যেষ্ঠ আশারিদু থেকে।

রাজা এসারহাদ্দোনকে মার-ইশতারের লেখা একটি চিঠিতে (খ্রিস্টপূর্ব ৬৬৮) আমরা পড়ি:

…প্রথম মাসে… ২৯ তারিখে বৃহস্পতিকে সরিয়ে নেওয়া হলো… এখন সে এক মাস পাঁচ দিন আকাশে অনুপস্থিত থেকেছে [অর্থাৎ অদৃশ্য ছিল]; তৃতীয় মাসের ষষ্ঠ দিনে বৃহস্পতি কালপুরুষ অঞ্চলে দৃশ্যমান হলো; তার নির্ধারিত সময়ের চেয়ে পাঁচ দিন বেশি সে পার করছিল। এটি এই লক্ষণের সঙ্গে মেলে: বৃহস্পতি তৃতীয় মাসে দেখা দিলে দেশ বিধ্বস্ত হবে এবং শস্য দুর্মূল্য হবে… বৃহস্পতি কালপুরুষে প্রবেশ করলে দেবতারা দেশকে গ্রাস করবেন।

এই এবং এ জাতীয় অন্যান্য লেখায় (যেখানে কিছু অভিব্যক্তির অর্থ এখনো অনিশ্চিত) আমরা এমন জ্যোতির্বৈজ্ঞানিক পর্যবেক্ষণ পাই, যা আগে পঞ্জিকা বা দিক নির্ণয়ের জন্য যা দরকার হতো তার চেয়ে অনেক বেশি বৈচিত্র্যময় ও বিস্তারিত। জ্যোতিষশাস্ত্র মানুষের জীবনকে আকাশের সঙ্গে এত ঘনিষ্ঠভাবে যুক্ত করেছিল যে নক্ষত্র ও তাদের পরিভ্রমণ তার চিন্তা ও কর্মকাণ্ডে গুরুত্বপূর্ণ স্থান নিতে শুরু করল। উজ্জ্বল দেবতারা সেই আশ্চর্য কক্ষপথগুলোতে তার ভাগ্যই বুনে চলতেন। এখন যেহেতু তার চোখ খুলে গিয়েছিল এবং আগ্রহ জেগেছিল, মানুষ গ্রহগুলোর গতিপথের সঙ্গে ক্রমশ আরও পরিচিত হয়ে উঠল। প্রতিবেদনগুলোতে আমরা তা দেখেছি: বৃহস্পতি পশ্চিমে যায়; মঙ্গল বৃশ্চিকে দাঁড়িয়ে ছিল, ঘুরে কম উজ্জ্বলতা নিয়ে এগিয়ে যায়; শুক্র তার জায়গায় থাকে; আর শনি সিংহে দেখা দিয়েছে। পর্যায়কালের কোনো তত্ত্ব ছিল না (এই সব বিবৃতিতে বছর ও তারিখ অনুপস্থিত)। তবু পর্যবেক্ষকদের মনে নিশ্চয়ই একটি পরিচিত ও প্রত্যাশিত নিয়মিততার ধারণা জন্মেছিল। গ্রহটি সঠিক জায়গায় ছিল বা তার সময় পেরিয়ে গেছে, এমন বিবৃতিতে তা প্রকাশ পায়।

গ্রহের পরেই জ্যোতিষীর কাছে প্রধান জ্যোতিষ্ক ছিল চাঁদ। এর কালগণনাগত তাৎপর্য এখন জ্যোতিষীয় আগ্রহের কাছে জায়গা ছেড়ে দিতে বাধ্য হলো। কালগণনার ত্রুটিগুলোও জ্যোতিষীদের কাছে ভবিষ্যদ্বাণীর উপাদান হয়ে উঠল। চাঁদ সময়মতো না এলে তা বেশিরভাগ ক্ষেত্রে ছিল অশুভ লক্ষণ। মেঘের কারণে চাঁদ অদৃশ্য থাকায়, কিংবা হয়তো দুর্যোগের সময়ে অবহেলার কারণে, নতুন বাঁকা চাঁদ মাসের ২৮ বা ২৯ তারিখেও দেখা দিতে পারত। একইভাবে পূর্ণিমা, যা স্বাভাবিকভাবে ১৪ তারিখে হওয়ার কথা, ১৩, ১৫ বা ১৬ তারিখে হতে পারত।

চাঁদ প্রথম দিনে দেখা দিলে নীরবতা থাকবে, দেশ পরিতৃপ্ত থাকবে।…

নিসান্নুর ত্রিশ তারিখে চাঁদ দেখা দিলে সুবারতু আখলামুকে গ্রাস করবে; আমুররুতে একটি বিদেশি ভাষা প্রাধান্য পাবে। আমরাই সুবারতু। ত্রিশতম দিনে চাঁদ দেখা দিলে দেশে শীত পড়বে। তেবিতুর ১৪ তারিখে চাঁদ সূর্য ছাড়া দেখা দিয়েছে; শাবাতুতে চাঁদ দিন পূর্ণ করবে।…

ত্রয়োদশ দিনে চাঁদ ও সূর্যকে একসঙ্গে দেখা গেলে নীরবতা থাকবে না; দেশে ব্যবসা-বাণিজ্য ব্যর্থ হবে, শত্রু দেশ দখল করবে। আপলা থেকে।

চাঁদ সূর্যের কাছে পৌঁছে তার সঙ্গে দৃষ্টির আড়ালে মিলিয়ে গেলে… দেশে সত্য প্রতিষ্ঠিত হবে এবং পুত্র পিতার সঙ্গে সত্য কথা বলবে। ১৪ তারিখে দেবতাকে দেবতার সঙ্গে দেখা গেছে।… ১৪ তারিখে চাঁদ ও সূর্যকে একসঙ্গে দেখা গেলে নীরবতা থাকবে, দেশ পরিতৃপ্ত থাকবে; দেবতারা আক্কাদের সুখ চান।…

চাঁদ সূর্যের জন্য অপেক্ষা না করে অদৃশ্য হলে সিংহ ও নেকড়ের তাণ্ডব হবে।… ১৫ তারিখে একে সূর্যের সঙ্গে দেখা গেছে; পরে তিশরিতুতে চাঁদ দিন পূর্ণ করবে।… বালাসি থেকে।

প্রথম দিনে আমি রাজার কাছে এভাবে পাঠিয়েছিলাম: ১৪ তারিখে চাঁদকে সূর্যের সঙ্গে দেখা যাবে।… ১৪ তারিখে চাঁদকে সূর্যের সঙ্গে দেখা গেছে।

ষোড়শ দিনে চাঁদ ও সূর্যকে একসঙ্গে দেখা গেলে রাজা রাজার প্রতি শত্রুতা পাঠাবেন। রাজা এক মাস ধরে তাঁর প্রাসাদে অবরুদ্ধ থাকবেন। শত্রুর পা দেশের বিরুদ্ধে থাকবে; শত্রু বিজয়গর্বে তার দেশে কুচকাওয়াজ করবে। দুজুর ১৪ বা ১৫ তারিখে চাঁদকে সূর্যের সঙ্গে দেখা না গেলে রাজা তাঁর প্রাসাদে অবরুদ্ধ হবেন। ষোড়শ দিনে দেখা গেলে তা সুবারতুর জন্য শুভ, আক্কাদ ও আমুররুর জন্য অশুভ। আকেল্লানু থেকে।

এমন পর্যবেক্ষণ ও সিদ্ধান্ত বুঝতে হলে পূর্ণিমার কাছাকাছি সময়ে কোন ঘটনাগুলো দেখা যায় তা বিবেচনা করতে হবে। মাস যদি স্বাভাবিকভাবে শুরু হয়, অর্থাৎ যে সন্ধ্যায় বাঁকা চাঁদ প্রথমবার দৃশ্যমান হয় (সাধারণত সংযোগের ১.৪ দিন পর) সেদিন থেকেই প্রথম দিন শুরু হয়, তাহলে পূর্ণিমা পড়বে ১৪ তারিখের রাতে। কারণ পূর্ণিমা আসে সংযোগের ১৪.৭ দিন পর (গড়ে; চাঁদের গতির অনিয়মের কারণে এটি ½ দিন কম বা বেশি হতে পারে)। ১৩ তারিখের রাতে চাঁদ তখনো পূর্ণ হয়নি। সূর্যাস্তের সময়েই চাঁদ পূর্বে দৃশ্যমান থাকে, আর সূর্যোদয়ের আগেই সে অস্ত যায়; সে সূর্যের জন্য অপেক্ষা করে না। ১৪ তারিখের রাতেও সূর্যাস্তের সময় চাঁদ দৃশ্যমান থাকে, কিন্তু সূর্য ওঠার সময়ও সে অস্ত যায়নি, পশ্চিমে নিষ্প্রভ ও ম্লান হয়ে আছে। আসিরীয় জ্যোতিষীরা একে প্রকাশ করতেন এভাবে: "চাঁদকে সূর্যের সঙ্গে দেখা গেছে", বা "দেবতা [সিন] দেবতা [শামাশ]-এর সঙ্গে দেখা দিয়েছেন", বা "চাঁদ সূর্যের কাছে পৌঁছায়"। এমন স্বাভাবিক অবস্থার সঙ্গে নানা রকম শুভ লক্ষণ জড়িত ছিল। কিন্তু যদি "চাঁদ সূর্যের জন্য অপেক্ষা না করে" সূর্যোদয়ের আগেই অস্ত যায়, তাহলে পূর্ণিমা আসবে পরে। তখন পরের রাত, অর্থাৎ ১৫ তারিখের আগে প্রভাতে চাঁদকে সূর্যের সঙ্গে দেখা যাবে না।

এই পর্যবেক্ষণগুলোতে জ্যোতিষশাস্ত্র কালগণনার সঙ্গে মিশে গিয়েছিল। পুরোহিত-জ্যোতির্বিদেরা লক্ষ করতেন যে আকাশে সব ঠিকঠাক আছে, তাই পৃথিবীতে শান্তি থাকবে। অথবা তাঁরা পঞ্জিকায় কোনো বিশৃঙ্খলা লক্ষ করতেন, যা অশুভ লক্ষণ ছাড়াও নির্দেশ করত যে কিছু একটা সংশোধন করতে হবে। প্রথম আবির্ভাবের দিন থেকে হিসাব করলে পূর্ণিমা ১৪ তারিখে প্রত্যাশিত ছিল, কিন্তু অনিয়মের কারণে তা না-ও হতে পারত। পূর্ণিমা ১৩ তারিখে এলে তাঁরা সিদ্ধান্ত নিতেন যে মাসটি হয়তো মাত্র ২৯ দিনের হবে। ১৪ তারিখের পরে চাঁদকে সূর্যের সঙ্গে দেখা গেলে মাসটিকে ৩০ দিনের হতে হতো; তখন সে "দিন পূর্ণ করত"। ১৪ তারিখে একসঙ্গে দেখা গেলে এবং আগের মাস ৩০ দিনের হলে পরের মাসটি হতো মাত্র ২৯ দিনের।

আদারুতে চাঁদ দিন পূর্ণ করেছে; ১৪ তারিখে চাঁদকে সূর্যের সঙ্গে দেখা যাবে; নিসান্নুতে চাঁদ দিন ফিরিয়ে নেবে।…

এভাবে জ্যোতিষীয় উদ্দেশ্য চাঁদ পর্যবেক্ষণে বেশি মনোযোগ দাবি করে কালগণনার জন্য আরও ভালো ভিত্তি জুগিয়েছিল। শৃঙ্খলা ও নিয়মিততা উপলব্ধি করা হলো এবং ভবিষ্যদ্বাণীর কাজে লাগানো হলো।

অবশ্যই চাঁদে আরও অনেক ঘটনা পর্যবেক্ষণ করা যেত: এর রং ও উজ্জ্বলতা, শিংয়ের আকৃতি, পৃথিবী-আলো (চাঁদ একটি "আগু", অর্থাৎ মুকুট বা রাজকীয় টুপি পরে আছে), একটি কুয়াশাময় ছটা, বা একটি বলয় (জ্যোতির্বলয়)। এমন বলয়কে প্রায়ই মনে করা হতো ভেড়ার চারণভূমি ঘিরে রাখা একটি বেড়া, যার কেন্দ্রে রাখাল হিসেবে আছে চাঁদ। গ্রহের ব্যাবিলনীয় চিহ্ন "লুবাত"-এর অর্থ "পথহারা ভেড়া"। অথবা বলয়টিকে একটি নদী হিসেবে দেখা হতো; কিংবা এর অর্থ ছিল অবরোধ, আর বলয়ের ভেতরের গ্রহ বা নক্ষত্রগুলো নির্দেশ করত কে অবরুদ্ধ। বলয়টি বন্ধ না হলে স্বাভাবিকভাবেই তার অর্থ ছিল পলায়ন।

জ্যোতির্বলয় চাঁদকে ঘিরে থাকলে এবং বৃহস্পতি তার ভেতরে থাকলে রাজা অবরুদ্ধ হবেন। বলয়টি বিচ্ছিন্ন ছিল; এটি অশুভ নির্দেশ করে না।… নাবু-শুমা-ইশকুন থেকে।

জ্যোতির্বলয় চাঁদকে ঘিরে থাকলে এবং সুদুন তার ভেতরে থাকলে একজন রাজা মারা যাবেন এবং তাঁর দেশ সংকুচিত হবে; এলামের রাজা মারা যাবেন। সুদুন হলো মঙ্গল, মঙ্গল আমুররুর নক্ষত্র; এটি আমুররু ও এলামের জন্য অশুভ। শনি আক্কাদের নক্ষত্র, এটি রাজা, আমার প্রভুর জন্য শুভ। রাজার সেবক ইরাশ্‌শি-ইলু থেকে।

গ্রহণকেও অত্যন্ত গুরুত্বপূর্ণ লক্ষণ মনে করা হতো। সব যুগেই এর অপ্রত্যাশিত আবির্ভাব মানুষকে গভীরভাবে নাড়া দিয়েছে। নিয়মিত ঘটনাগুলোকেই যেখানে এত তাৎপর্য দেওয়া হতো, সেখানে এর প্রভাব আরও বেশি ছিল। গ্রহণ থেকে বহু রকম ভবিষ্যদ্বাণী করা হতো। সূর্যগ্রহণের ক্ষেত্রে মাস ও আকাশে অবস্থান লিখে রাখা হতো। সূর্য যখন "চাঁদের আকার ধারণ করত", তখন শিংগুলোর চেহারা ও দিকও লিখে রাখা হতো। মাসের ২৭ বা ২৮ তারিখ ছাড়া অন্য তারিখে গ্রহণের বিবরণ পাওয়া গেলে বোঝা যায়, জ্যোতির্বৈজ্ঞানিক গ্রহণ ছাড়াও ধুলিঝড় সূর্যকে অন্ধকার করে থাকতে পারে।

চন্দ্রগ্রহণ থেকে পাওয়া লক্ষণ ছিল অনেক বেশি। চাঁদের পূর্ণগ্রহণ হলে তা সব দেশের জন্য লক্ষণ ছিল। আংশিক আচ্ছন্ন হলে এর চার দিকের প্রতিটি আলাদা একটি দেশের সঙ্গে সম্পর্কিত ছিল, আর মাস, দিন ও প্রহর (রাতের পাহারা) অনুযায়ী তাৎপর্য আলাদা হতো। তাই "এনুমা আনু এনলিল" নামের প্রাচীন লক্ষণের বিশাল সংকলনে এই নির্দেশনা দেওয়া হয়েছে:

চাঁদের গ্রহণ হলে তুমি নিখুঁতভাবে পর্যবেক্ষণ করবে মাস, দিন, রাতের প্রহর, বাতাস, গতিপথ এবং সেই নক্ষত্রগুলোর অবস্থান যাদের এলাকায় গ্রহণ ঘটছে। এর মাস, দিন, রাতের প্রহর, বাতাস, গতিপথ ও নক্ষত্র সম্পর্কিত লক্ষণগুলো তুমি নির্দেশ করবে।

তাই সিমান্নু মাসে রাতের শেষে ঘটে যাওয়া একটি গ্রহণের প্রতিবেদনে আমরা পাই:

প্রভাতী প্রহরে গ্রহণের অর্থ রোগ।… প্রভাতী প্রহর হলো এলাম, ১৪ তারিখ হলো এলাম, সিমান্নু হলো আমুররু, দ্বিতীয় দিক হলো আক্কাদ।… প্রভাতী প্রহরে গ্রহণ হয়ে প্রহরটি পূর্ণ করলে, এবং উত্তুরে বাতাস বইলে, আক্কাদের অসুস্থরা সেরে উঠবে। গ্রহণ প্রথম দিকে শুরু হয়ে দ্বিতীয় দিকে থামলে এলামে হত্যাযজ্ঞ হবে; গুতি আক্কাদের কাছে আসবে না।… গ্রহণ হয়ে দ্বিতীয় দিকে থামলে দেবতারা দেশের প্রতি দয়া দেখাবেন। সিমান্নুতে চাঁদ অন্ধকার হলে এক বছর পর রামানু [ঝড়ের দেবতা] প্লাবন ঘটাবেন। সিমান্নুতে চাঁদের গ্রহণ হলে বন্যা হবে, আর দেশের জলজ উৎপাদন প্রচুর হবে।…

লক্ষণের এক সমৃদ্ধ তালিকা; প্রতিটি খুঁটিনাটির তাৎপর্য ছিল। এত বিস্তারিত বিবরণ থাকলে আমরা গ্রহণের বেশ নির্ভুল পর্যবেক্ষণের কথা বলতে পারি। তাই আশ্চর্যের কিছু নেই যে বহু শতাব্দী পরে টলেমি চাঁদের পর্যায়কাল নির্ণয়ের জন্য যে প্রথম পর্যবেক্ষিত গ্রহণগুলো ব্যবহার করেছিলেন, সেগুলো এই যুগেরই: খ্রিস্টপূর্ব ৭২১ সালের ১৯ মার্চ, খ্রিস্টপূর্ব ৭২০ সালের ৮ মার্চ এবং খ্রিস্টপূর্ব ৭২০ সালের ১ সেপ্টেম্বর। এখানে কেবল প্রয়োজন ছিল এই তারিখগুলোর পর কত বছর ও কত দিন পেরিয়েছে তার নিখুঁত জ্ঞান। হতে পারে, এই তথ্যের অভাবেই আরও আগের গ্রহণগুলো ব্যবহার করা যায়নি। আর এ কারণেই টলেমি তাঁর রাজাদের তালিকা ও তাঁদের শাসনের ধারাবাহিক বছরগুলো শুরু করেছিলেন ব্যাবিলনীয় রাজা নাবোনাসারকে (খ্রিস্টপূর্ব ৭৪৭) দিয়ে।

ব্যাবিলনীয় প্রহরীরা নিজেরা নিশ্চয়ই এত যত্নের সঙ্গে পর্যবেক্ষিত ঘটনাবলিতে নিয়মিততা লক্ষ করেছিলেন। তাঁদের প্রতিবেদন দেখায় যে তাঁরা গ্রহণ প্রত্যাশা করতেন এবং ঘটনাগুলো আগাম ঘোষণা করতেন:

১৪ তারিখে একটি গ্রহণ ঘটবে; এটি এলাম ও আমুররুর জন্য অশুভ, রাজা, আমার প্রভুর জন্য শুভ; রাজা, আমার প্রভু, নিশ্চিন্তে থাকুন। এটি শুক্র ছাড়াই দেখা যাবে। রাজা, আমার প্রভুকে আমি বলছি: একটি গ্রহণ হবে। রাজার সেবক ইরাশ্‌শি-ইলু থেকে।

দেশসমূহের রাজা, আমার প্রভুর প্রতি, তোমার সেবক বিল-উসুর। বেল, নেবো ও শামাশ রাজা, আমার প্রভুর প্রতি সদয় হোন। একটি গ্রহণ ঘটেছে, কিন্তু রাজধানীতে তা দেখা যায়নি। সেই গ্রহণ যখন এগিয়ে আসছিল, তখন রাজা যেখানে বাস করেন সেই রাজধানীতে, দেখুন, সর্বত্র মেঘ ছিল। গ্রহণ ঘটেছে কি ঘটেনি তা আমরা জানি না। রাজাদের প্রভু আশুরে, সব নগরে, ব্যাবিলন, নিপ্পুর, উরুক ও বোরসিপ্পায় লোক পাঠান; সেই নগরগুলোতে যা দেখা গেছে রাজা তা নিশ্চিতভাবে শুনবেন।… রাজা, আমার প্রভুর নগরে বসবাসকারী মহান দেবতারা আকাশ মেঘে ঢেকে দিয়েছেন এবং গ্রহণ দেখতে দেননি। তাই রাজা জানুন যে এই গ্রহণ রাজা, আমার প্রভু, বা তাঁর দেশের বিরুদ্ধে নয়। রাজা আনন্দ করুন।…

এটি কি চমৎকার একটি ভাবনা নয় যে গ্রহণ ঘটার সময় আশুরের রক্ষাকর্তা দেবতারা, দেশ ও রাজার কোনো বিপদ নেই জেনে, নগরের ওপর মেঘের পর্দা টেনে দিলেন, যাতে রাজা ভয় না পান? লাইডেনের অ্যাসিরিওলজিস্ট দে লিয়াগ্রে ব্যোলের মতে, এই প্রতিবেদনগুলোর বেশিরভাগই খ্রিস্টপূর্ব ৬৭৫ থেকে ৬৬৫ সালের মধ্যেকার। তখন ছিল আসিরীয় সাম্রাজ্যের সংকটকাল; অসুস্থ ও আশঙ্কাগ্রস্ত রাজা কমবেশি ব্যর্থ একটি মিশর অভিযান থেকে ফিরেছিলেন। এখানে আমরা একটি বেশ নির্ভুল তারিখ পাই, যে সময়ের জ্যোতির্বৈজ্ঞানিক জ্ঞানের অবস্থা বলা যায়। এটা সত্যিই জ্ঞান ছিল, কারণ এই প্রতিবেদনগুলোতে গ্রহণ সম্পর্কে ভবিষ্যদ্বাণী করা হয়েছে এবং খুব স্পষ্টভাবে তা যাচাইও করা হয়েছে।

রাজা, আমার প্রভুর কাছে আমি পাঠিয়েছিলাম: "একটি গ্রহণ ঘটবে।" এখন তা ব্যর্থ হয়নি, ঘটেছে। এই গ্রহণ ঘটায় এটি রাজা, আমার প্রভুর জন্য শান্তির পূর্বাভাস দেয়।…

পুরোহিতেরা কিসের ভিত্তিতে ভবিষ্যদ্বাণী করতেন, এই লেখাগুলোতে তার কোনো ইঙ্গিত নেই। তবে চন্দ্রগ্রহণ ঘটার মধ্যে এমন একটি সরল নিয়মিততা আছে (যা প্রথম দেখিয়েছিলেন স্কিয়াপারেল্লি) যে তা নিশ্চয়ই শিগগিরই তাঁদের নজরে পড়েছিল। তাঁরা দেখতে পেতেন, ছয় মাসের কম ব্যবধানে কখনো দ্বিতীয় গ্রহণ হয় না, আর প্রায়ই আরও ছয় মাস পর তৃতীয় একটি গ্রহণ হয়। কখনো কখনো এমনকি চার বা পাঁচটি গ্রহণ পরপর ঘটত, সবসময় ছয় মাসের ব্যবধানে। প্রায়ই এই ধারাবাহিক শ্রেণিতে একটি ফাঁক থাকত। এমন হতো যখন গ্রহণটি দিনের বেলায় ঘটত, ফলে পূর্ণিমার চাঁদ দিগন্তের নিচে থাকত। পর্যবেক্ষকেরা একবার যখন বুঝলেন যে এটাই ফাঁকের কারণ, তখন সঠিক মুহূর্তে একটি অদৃশ্য গ্রহণ যোগ করে তাঁরা ফাঁকটি পূরণ করতে পারতেন। এভাবে পাঁচ বা ছয়টি দৃশ্যমান বা অদৃশ্য গ্রহণের একটি নিয়মিত ও অবিচ্ছিন্ন শ্রেণি তৈরি হতো। তারপর শ্রেণিটি শেষ হয়ে যেত। কিন্তু এক বা দুই বছর পর একটি নতুন শ্রেণি শুরু হতো, আগের শ্রেণিটি চলতে থাকলে যে সময়ে হতো তার এক মাস আগে। তাই যদি তাঁরা এমন একটি গ্রহণ পর্যবেক্ষণ করতেন যার ৬, ১২ বা ১৮ মাস আগে আর কোনো গ্রহণ হয়নি, তাহলে সিদ্ধান্ত নিতে পারতেন যে একটি নতুন শ্রেণি শুরু হয়েছে। এভাবে ৬ বা ১২ মাস পরের নতুন গ্রহণের ভবিষ্যদ্বাণী করতে পারতেন।

আধুনিক জ্ঞানের আলোকে এই নিয়মিততার ভিত্তি সহজেই বোঝা যায়। ছয়টি চান্দ্র পর্যায়কাল (১৭৭.১৮ দিন) পর সূর্য, আর তাই পূর্ণিমার চাঁদের অবস্থানও, দ্রাঘিমায় গড়ে ১৭৪.৬৪° এগিয়ে যায়। সূর্য ও চাঁদের পরিবর্তনশীল বেগের কারণে প্রকৃত অগ্রগতি কয়েক ডিগ্রি কম বা বেশি হতে পারে। একই সময়ে চন্দ্রকক্ষের পাতবিন্দুগুলো (nodes) ৯.৩৮° পিছিয়ে যায়; তাই বিপরীত পাতবিন্দুটি ১৭০.৬২° বেশি দ্রাঘিমায় অবস্থান করে। পাতবিন্দুর সাপেক্ষে পূর্ণিমার চাঁদ তখন ৪.০২° এগিয়ে যায়। চাঁদ যদি কোনো পাতবিন্দুর কাছে থাকে, তাহলে ছয় মাস পর আবারও একটি পাতবিন্দুর কাছে থাকবে। পাতবিন্দু থেকে পূর্ণিমার চাঁদের দূরত্ব ১০° থেকে ১২°-এর বেশি না হলে চাঁদ পৃথিবীর ছায়ার মধ্য দিয়ে যেতে পারে। পূর্ণগ্রহণের জন্য দূরত্ব ৫° বা ৬°-এর বেশি হওয়া চলবে না, নইলে গ্রহণটি কেবল আংশিক হয়। ছয় মাসের ব্যবধানে পরপর কয়েকটি পূর্ণিমার দূরত্ব হিসাব করলে এই শর্তগুলো স্পষ্ট করা যায়। হিসাবটি যেকোনো একটি প্রথম মান থেকে শুরু করে প্রতিবার ৪.০২ বা ৪.০৩ করে বাড়িয়ে করা হয়:

দূরত্ব	−১৪.৫০	−১০.৪৮	−৬.৪৫	−২.৪৩	+১.৫৯	+৫.৬২	+৯.৬৫	+১৩.৬৭
ফল	গ্রহণ নেই	আংশিক?	আংশিক	পূর্ণ	পূর্ণ	পূর্ণ?	আংশিক	গ্রহণ নেই

সংশ্লিষ্ট চিত্রে (২) এই অবস্থাগুলো দেখানো হয়েছে। সেখানে বড় বৃত্তগুলো ছায়া এবং ছোট বৃত্তগুলো চাঁদকে বোঝায়।